\HMTNarrativeHeading{Persistence and electronic trajectory}

The electronic structure begins at the persistent central position of APP.
Cellular inversion distinguishes this position from the conjugate pairs
that surround it. Its evaluation preserves the sheets, remainders, and
quotients; the trajectory prolongs these data through the operations of the
TPK. The exposition of mass subsequently receives the action sections and
the readers it uses.

The central vacancy makes visible a decisive difference between residual
reading and arithmetic history. Two positions may share the sum and the
multiplicative remainder, while their integer products preserve distinct
quotients. The change in quotient records the transfer; orientation
distinguishes their conjugate positions. The example makes visible the
information needed by a reconstruction that restores the operation
performed.

The electronic chain adds an explicit chronology. The cellular recurrence,
inversions, and reading windows preserve the registers of the persistent
section. The calendar records the steps at which capacity remains
unchanged; the central defect records the difference in quotient between
evaluations with the same remainder. The electronic trajectory preserves
this second datum together with its orientation and chronology. Capacity
vacancy and quotient defect count distinct operations, with their
respective registers.

This continuity makes the distinction between return and recovery
pertinent once again. A phase may recur while the register of the section
incorporates the path. The subsequent evaluation uses this transported
information and keeps the central section identified. Mass, orientation,
and the expressions of action are presented together with their
antecedents, so that each step can be read as a definite composition.

The electron thus provides a developed case of the article's thesis.
Persistence corresponds to a structure; the trajectory describes its
evolution; the registers distinguish the changes; the dimensional reader
obtains a magnitude. Its relevance to conservation lies in this complete
sequence. The unity maintained during transport preserves relations whose
content becomes accessible through different readings. The subsequent
operator study will abstract the conditions for recovery and show their
exact balances.
